% =====================================================================
%  3. Delineamento metodológico
% =====================================================================
\section{Delineamento metodológico}
\label{sec:delineamento}

\subsection{Premissas}

O encadeamento das atividades parte de duas premissas estabelecidas no MOP
\cite{mop2026}:

\begin{enumerate}[label=\alph*)]
  \item as comunidades contempladas serão definidas por meio do trabalho
        técnico-social (\autoref{subsec:tts});
  \item o modelo de gestão será definido e pactuado com as comunidades
        selecionadas \textbf{antes} da instalação dos sistemas
        (\autoref{subsec:modelos-gestao}).
\end{enumerate}

\subsection{Estrutura das atividades}

O Subcomponente 3.2.b é composto por seis atividades, agrupadas em três
blocos sequenciais (\autoref{fig:fluxograma-geral}):

\begin{enumerate}[label=\roman*)]
  \item \textbf{estruturação institucional e social}: implantação dos
        modelos de gestão (\atividade{3.2.2.1}) e trabalho técnico-social
        (\atividade{3.2.2.2});
  \item \textbf{implantação física}: sistemas coletivos de abastecimento de
        água (\atividade{3.2.2.3}), sistemas descentralizados de esgotamento
        sanitário em comunidades rurais (\atividade{3.2.2.4}) e na área de
        atuação do IDR-Paraná (\atividade{3.2.2.5}), e sistemas do Programa
        Sanepar Rural (\atividade{3.2.2.6});
  \item \textbf{entrega, monitoramento e encerramento}: transferência dos
        sistemas às entidades gestoras, capacitação dos beneficiários,
        avaliação de desempenho e prestação de contas.
\end{enumerate}

As salvaguardas ambientais e sociais (\autoref{sec:salvaguardas}) permeiam
todos os blocos.

\begin{figure}[htbp]
  \centering
  \caption{Fluxograma metodológico do Subcomponente 3.2.b}
  \label{fig:fluxograma-geral}
  % Fluxograma metodológico geral do Subcomponente 3.2.b
\begin{tikzpicture}[node distance=0.55cm and 0.4cm]

  % --- Bloco I: estruturação institucional e social -----------------
  \node[marco] (sel) {Seleção de áreas e comunidades\\
    \normalfont\scriptsize critérios estaduais de seleção};

  \node[etapa=7.4cm, below=of sel] (gestao) {\textbf{3.2.2.1 -- Modelos de gestão}\\
    definição do modelo, pactuação e termo de cooperação};

  \node[etapa=7.4cm, below=of gestao] (tts) {\textbf{3.2.2.2 -- Trabalho técnico-social}\\
    mobilização, diagnóstico, enquadramento e adesão das famílias};

  \node[externo, left=0.5cm of tts] (concessao) {Municípios com contrato de concessão com a Sanepar};
  \node[externo, right=0.5cm of tts] (pip) {PIP elaborados pelo IDR-Paraná (Subcomp.~2.3)};

  % --- Bloco II: implantação física --------------------------------
  \node[etapa, below=1.1cm of concessao] (sanepar) {\textbf{3.2.2.6}\\ Sanepar Rural\\
    \scriptsize convênio -- contrapartida};
  \node[etapa, below=1.1cm of pip] (esgidr) {\textbf{3.2.2.5}\\ Esgotamento -- área IDR\\
    \scriptsize 2.300 sistemas};

  \path (sanepar) -- (esgidr) coordinate[pos=0.5] (meio);
  \node[etapa=3.2cm] (agua) at ($(meio)+(-2.0cm,0)$) {\textbf{3.2.2.3}\\ Abastecimento de água\\
    \scriptsize 83 comunidades};
  \node[etapa=3.2cm] (esgcom) at ($(meio)+(2.0cm,0)$) {\textbf{3.2.2.4}\\ Esgotamento -- comunidades\\
    \scriptsize 4.067 domicílios};

  % --- Bloco III: entrega, monitoramento e encerramento ------------
  \node[etapa=15.2cm, anchor=north] (entrega) at ($(meio |- agua.south)+(0,-1.1cm)$) {\textbf{Entrega e capacitação}\\
    capacitação dos beneficiários, entrega à entidade gestora ou ao proprietário e aceite formal};

  \node[marco=15.2cm, below=of entrega] (fim) {Monitoramento, avaliação e encerramento\\
    \normalfont\scriptsize relatórios semestrais, indicadores, relatório final e Termo de Recebimento Definitivo};

  \node[transversal=15.2cm, below=0.45cm of fim] (salv) {\textbf{Transversal:}
    salvaguardas ambientais e sociais -- MGAS, PEPI, PGR, MRF e MQR (\autoref{sec:salvaguardas})};

  % --- Ligações ----------------------------------------------------
  \draw[seta] (sel) -- (gestao);
  \draw[seta] (gestao) -- node[rotulo, right=2pt] {modelo pactuado antes das obras} (tts);
  \draw[seta] (tts.south) ++(-2.0cm,0) -- (agua.north);
  \draw[seta] (tts.south) ++(2.0cm,0) -- (esgcom.north);
  \draw[seta] (agua) -- node[rotulo, above] {70\%} (esgcom);
  \draw[seta] (concessao) -- (sanepar);
  \draw[seta] (pip) -- (esgidr);
  \draw[setatracejada] (gestao.west) -| (sanepar.north west);

  \foreach \n in {sanepar, agua, esgcom, esgidr}
    \draw[seta] (\n.south) -- (\n.south |- entrega.north);
  \draw[seta] (entrega) -- (fim);

  % --- Rótulos dos blocos -------------------------------------------
  \begin{scope}[on background layer]
    \node[fill=cinzaclaro!60, rounded corners, inner sep=6pt,
          fit=(sel)(gestao)(tts)(concessao)(pip)] (b1) {};
    \node[fill=cinzaclaro!60, rounded corners, inner sep=6pt,
          fit=(sanepar)(agua)(esgcom)(esgidr)] (b2) {};
    \node[fill=cinzaclaro!60, rounded corners, inner sep=6pt,
          fit=(entrega)(fim)] (b3) {};
  \end{scope}
  \node[font=\scriptsize\bfseries, anchor=south west, text=azulpsh]
    at (b1.north west) {I. Estruturação institucional e social};
  % rótulos centrados no espaço livre entre as setas de 3.2.2.6 e 3.2.2.3
  \path (sanepar.center) -- (agua.center) coordinate[pos=0.5] (livre);
  \node[font=\scriptsize\bfseries, anchor=south, text=azulpsh]
    at (livre |- b2.north) {II. Implantação física};
  \node[font=\scriptsize\bfseries, anchor=south, text=azulpsh]
    at (livre |- b3.north) {III. Entrega e encerramento};
\end{tikzpicture}

  \fonte{Elaborado pelo IAT (2026) com base em \citeonline{mop2026}.}
\end{figure}

\subsection{Arranjo institucional e financiamento}
\label{subsec:arranjo}

A ação será executada pelo Instituto Água e Terra (IAT), com coexecução do
Instituto de Desenvolvimento Rural do Paraná (IDR-Paraná) e da Companhia de
Saneamento do Paraná (Sanepar). As \atividade{3.2.2.1} a \atividade{3.2.2.5}
serão executadas por contratação pública, mediante licitação, com
financiamento do Banco Mundial (BIRD), em um período de seis anos e custo
total estimado de US\$ 35,48 milhões. A Sanepar participará com contrapartida
estimada em US\$ 12,31 milhões, destinada ao Programa Sanepar Rural
(\atividade{3.2.2.6}) \cite{mop2026}. Os papéis de cada ator estão resumidos
no \autoref{qua:atores}.

\begin{quadro}[htbp]
\caption{Atores e papéis na execução do Subcomponente 3.2.b}
\label{qua:atores}
\small
\renewcommand{\arraystretch}{1.2}
\begin{tabularx}{\textwidth}{|P{3.6cm}|L|}
\hline
\cab{Ator} & \cab{Papel} \\ \hline
IAT & Executor. Planejamento, coordenação técnica, contratações, fiscalização,
      monitoramento e validação dos produtos. \\ \hline
IDR-Paraná & Coexecutor. Apoio ao planejamento, mobilização social,
      elaboração dos PIP e atuação nas áreas de mananciais
      (\atividade{3.2.2.5}). \\ \hline
Sanepar & Coexecutora. Execução do Programa Sanepar Rural com recursos de
      contrapartida (\atividade{3.2.2.6}). \\ \hline
SECID & Apoio ao planejamento e à coordenação técnica; validação da
      implantação dos modelos de gestão. \\ \hline
Prefeituras municipais & Participação na governança, articulação local e
      apoio às comunidades. \\ \hline
Comunidades e associações & Participação na gestão e na operação cotidiana
      dos sistemas implantados. \\ \hline
UGP / Banco Mundial & Consolidação dos relatórios, validação do
      monitoramento e do encerramento. \\ \hline
\end{tabularx}
\fonte{Adaptado de \citeonline{mop2026}.}
\end{quadro}

\pendencia{A forma de execução ``contratação pública'' está declarada para a
ação como um todo, mas o Programa Sanepar Rural (\atividade{3.2.2.6}) é
executado por convênio. O texto acima já faz essa distinção; confirmar com a
UGP.}
